%% main.tex — transmon-benchmark.1 (anvil:pub draft, iteration 1)
%%
%% Drafted 2026-07-14 from papers/transmon-benchmark/BRIEF.md (authoritative)
%% and the transmon-benchmark.0.litsearch sibling. All numbers trace to the
%% brief's tables and to committed artifacts (benchmarks/transmon_eigen/
%% results.toml, reference/fixtures/transmon_palace/results_p1/eig.csv).
%% GPU-cell numbers are PENDING: every would-be number is a \TBDGPU marker.
%% Do NOT replace a \TBDGPU marker with a number that is not in a committed
%% artifact.

\documentclass{anvil-paper}

%% --- URL line-breaking: the two AWS-blog references carry very long
%% hyphenated URLs; allow breaks at hyphens and add stretch so the
%% bibliography stays inside the measure (render-gate overfull check).
\expandafter\def\expandafter\UrlBreaks\expandafter{\UrlBreaks\do\-}
\Urlmuskip=0mu plus 1mu

%% --- Draft-phase helper macros -------------------------------------------
%% Unmistakable placeholder for pending GPU-cell content (brief: "TBD-GPU
%% placeholders must be impossible to mistake for results").
\newcommand{\TBDGPU}[1]{\textcolor{red!75!black}{\textbf{[TBD-GPU: #1]}}}
%% Figure placeholder: uses the rendered figure if pub-figures has produced
%% it, otherwise a framed placeholder naming the source script.
\newcommand{\anvilfig}[3][0.85]{%
  \IfFileExists{figures/#2}%
    {\includegraphics[width=#1\linewidth]{figures/#2}}%
    {\fbox{\parbox{#1\linewidth}{\centering\vspace{5em}%
      \texttt{figures/#2}\\[0.5em]%
      {\small placeholder --- render via \texttt{figures/src/#3}}%
      \vspace{5em}}}}}

\title{Cross-validating an open-source Rust FEM solver against Palace on a
superconducting transmon geometry}
\author{Robb Walters}
% TODO(operator): final author list and affiliation are operator-decided
% before submission (see BRIEF.md frontmatter).
\date{\today}

\begin{document}
\maketitle

\begin{abstract}
Palace, the open-source finite-element electromagnetics solver from the AWS
Center for Quantum Computing, has become a reference tool for
superconducting-qubit design, yet no independent quantitative validation of
it exists in the literature. We supply one. GEODE-FEM, an independently
implemented open-source Rust finite-element solver, reproduces Palace's
eigenmode solution on a real transmon-plus-readout-resonator geometry ---
generated by DeviceLayout.jl's \texttt{SingleTransmon} example and meshed
once, so both solvers consume the identical mesh and the identical
lumped-element junction model --- to within $0.03\%$ across all physical
modes, with the junction LC mode agreeing to $0.000\%$. Exact analytic
tripwires (the LC anchor and the Josephson $1/\sqrt{2}$ inductance scaling)
guard against coincidental agreement. On an 8-core CPU node, geode-fem's
serial direct-factorization shift-invert eigensolve is roughly $4\times$
more efficient per core than Palace's MPI-distributed Krylov solve at this
133k-DOF size, while Palace's parallelism wins $1.7\times$ at the whole-node
level --- an architecture trade-off, not a ranking. Two caveats are stated
up front: the model contains no ${\sim}4$\,GHz qubit mode \emph{by
construction} (both solvers agree on its absence), and geode-fem's
un-projected Lanczos admits one disclosed spurious junction-localized mode
that Palace's divergence-free projection suppresses. The full
fixture-to-comparison pipeline is scripted and committed.
\end{abstract}

\section{Introduction}
\label{sec:intro}

Three-dimensional finite-element electromagnetics is now a load-bearing step
in superconducting-qubit design: mode frequencies, dispersive shifts, and
participation ratios extracted from field simulation feed directly into
device Hamiltonians \citep{koch2007,nigg2012,minev2021}. The open-source
solver of record for this workflow is Palace \citep{palace}, an MFEM-based
\citep{mfem,andrej2024high} parallel solver from the AWS Center for Quantum
Computing, typically driven by geometry from DeviceLayout.jl
\citep{devicelayout,devicelayout-blog}. This stack is widely used and
well engineered --- and, unusually for scientific software of its
influence, has never been the subject of a peer-reviewed or preprint
publication: Palace exists as a repository, a blog post
\citep{palace-blog}, and a workshop talk; the DeviceLayout.jl transmon
workflow exists as a blog post and a conference talk
\citep{devicelayout-blog,devicelayout-juliacon}. Independent, quantitative,
third-party validation of the stack does not exist in the literature.

This paper supplies that validation, and does so under conditions designed
to make agreement meaningful rather than merely reassuring. We compare
Palace against GEODE-FEM (``geode-fem''), an open-source finite-element
solver implemented independently in Rust --- different language, different
assembly path, different eigensolver, different linear algebra --- on a real
transmon-plus-readout-resonator geometry produced by DeviceLayout.jl's own
\texttt{SingleTransmon} example. The geometry is meshed \emph{once}; both
solvers consume the identical mesh, the identical first-order discretization,
and the identical lumped reactive-shunt junction model. Agreement claims are
therefore about formulation and solver correctness, not meshing luck.

Our contributions are:
\begin{itemize}
  \item \textbf{The first independent validation of the Palace/DeviceLayout.jl
    transmon stack.} On the identical mesh at matched order, geode-fem
    reproduces Palace's physical eigenmodes to $\leq 0.03\%$, with the
    junction LC mode agreeing to $0.000\%$
    (Section~\ref{sec:agreement}).
  \item \textbf{A reproducible same-mesh cross-solver methodology.} An
    oracle-first benchmark protocol --- exact analytic tripwires, inverse
    tests that must fail, participation-based mode identification, and a
    fully scripted fixture-to-comparison pipeline
    (Sections~\ref{sec:methodology} and~\ref{sec:repro}).
  \item \textbf{An honest architecture-level performance comparison.} On an
    8-core node, a serial direct-factorization shift-invert eigensolve is
    ${\sim}4\times$ more efficient per core than a distributed Krylov solve
    at 133k DOFs, while MPI parallelism wins $1.7\times$ at the whole-box
    level (Section~\ref{sec:perf-cpu}); a GPU cell is scoped but pending
    (Section~\ref{sec:perf-gpu}).
  \item \textbf{Honest-negative reporting as a first-class result.} The
    model's lack of a ${\sim}4$\,GHz qubit mode is explained and confirmed
    by both solvers (Section~\ref{sec:absent-mode}); a spurious mode leaked
    by geode-fem's un-projected Lanczos is disclosed, characterized, and
    tied to a known fix (Section~\ref{sec:spurious-mode}).
\end{itemize}

Section~\ref{sec:related} positions this work; Section~\ref{sec:problem}
defines the benchmark problem; Section~\ref{sec:solvers} contrasts the two
solver architectures; the remaining sections present methodology, results,
and reproducibility details.

\section{Related Work}
\label{sec:related}

\paragraph{Electromagnetic simulation workflows for superconducting qubits.}
The closest prior work consists of two recent Palace-based workflow papers.
\citet{sommers2025open} present SQDMetal, an open workflow that couples
Qiskit Metal, Gmsh, and Palace, and cross-check Palace against the
commercial solvers COMSOL and Ansys HFSS on qubit and resonator geometries,
reporting agreement at roughly the $0.02$--$0.13\%$ level.
\citet{ye2025electromagnetic} build a Palace-centered layout-to-Hamiltonian
workflow and benchmark it against cryogenic measurement, with resonator
frequencies reproduced to within $0.3\%$. Both efforts treat Palace as a
component to be \emph{wrapped and trusted}: validation runs against
commercial tools or experiment, through independently constructed models
and meshes, so residuals mix solver error with meshing and model-entry
differences. Neither is a same-mesh, same-formulation comparison against an
independently \emph{implemented} solver. Our study eliminates meshing and
model entry as confounders --- identical MSH fixture, identical lumped-shunt
junction, first-order N\'ed\'elec elements on both sides --- so residuals
isolate formulation and solver correctness; to our knowledge it is the
first independent third-party validation of Palace itself. The broader
open ecosystem for qubit electromagnetic design includes Qiskit Metal and
its quasi-lumped analysis methods \citep{minev2021circuit} and the SQuADDS
validated design database \citep{shanto2024squadds}; the physics framing
for extracting qubit parameters from field modes rests on the transmon
\citep{koch2007}, black-box quantization \citep{nigg2012}, and
energy-participation-ratio \citep{minev2021} literature.

\paragraph{Verification, validation, and cross-solver benchmarking.}
Code-to-code comparison on precisely specified benchmark problems is an
established genre in computational electromagnetics: the TEAM (Testing
Electromagnetic Analysis Methods) workshop tradition has run such
comparisons since 1986.\footnote{The TEAM founding documents are workshop
reports and proceedings without resolvable DOIs; we describe the tradition
in prose rather than cite an arbitrary secondary source.} In the V\&V
vocabulary of \citet{roache1997quantification} and
\citet{oberkampf2010verification}, the present study is
verification-side evidence: our analytic tripwires are code-verification
instruments in exactly Roache's sense, and the same-mesh protocol is a
code-to-code comparison in Oberkampf and Roy's taxonomy. Multi-solver
comparisons on shared electromagnetic problems have precedent in
nano-optics, where \citet{hoffmann2009comparison} compared FEM, FDTD, and
FIT codes on plasmonic antennas. We import that discipline to the
superconducting-qubit domain and add two elements the genre often lacks:
exact analytic tripwires that must move in known ways, and honest-negative
reporting of what the model does \emph{not} contain.

\paragraph{Open-source FEM/EM solver ecosystems.}
Mature open electromagnetics tools are normally accompanied by a citable
methods paper: GetDP \citep{dular1998general}, DOLFIN/FEniCS
\citep{logg2010dolfin}, Gmsh \citep{gmsh}, and scikit-rf
\citep{arsenovic2022scikit} all have one; so does Palace's substrate
library MFEM \citep{mfem,andrej2024high}. Palace and DeviceLayout.jl are
outliers --- the \emph{library} is published, the solver product built on
it is not --- which is precisely the gap this benchmark narrows from the
outside, at a formality level (preprint) exceeding the original stack's
blog-and-talk record.

\paragraph{Numerics.}
geode-fem's eigenpath is a shift-invert Lanczos \citep{arpack} over a real
symmetric pencil, factorized with the faer sparse-LU \citep{kazdadi2026faer}
and assembled on the Burn tensor framework \citep{burn}; the element family
is first-order N\'ed\'elec \citep{nedelec1980}. Palace's eigenpath is a
SLEPc-class Krylov method \citep{slepc} with divergence-free projection and
auxiliary-space (AMS) preconditioning
\citep{hiptmair-xu,tzaniovkolevpanayotsvassilevski2018parallel}.

\section{Benchmark Problem: Geometry, Mesh, and Junction Model}
\label{sec:problem}

\subsection{Geometry and mesh provenance}
\label{sec:geometry}

The device is the \texttt{SingleTransmon} example shipped with
DeviceLayout.jl v1.15.0 \citep{devicelayout}: a transmon with readout
resonator on a sapphire substrate, enclosed in a vacuum box. The sapphire is
anisotropic, with permittivity tensor
$\varepsilon = R\,\mathrm{diag}(9.3,\, 9.3,\, 11.5)\,R^{\mathsf{T}}$
rotated approximately $36.87^\circ$ in-plane. The DeviceLayout.jl workflow
emits seven named physical groups (metal, substrate, vacuum, junction
lumped-element surface, two readout ports, exterior boundary) and drives
Gmsh \citep{gmsh} to produce a tetrahedral mesh, exported as an MSH~4.1
fixture and pinned by SHA-256 (Section~\ref{sec:repro}).

The mesh has 22{,}684 nodes and 133{,}314 tetrahedra, yielding 156{,}863
N\'ed\'elec edge unknowns, of which 133{,}108 interior degrees of freedom
remain after eliminating perfect-electric-conductor (PEC) boundary edges.
PEC is applied on all metal surfaces and the exterior box; the two readout
ports are left open (lossless) in this first version, so no resistive
element enters and the eigenproblem pencil stays real symmetric.
Figure~\ref{fig:geometry} shows the geometry with physical groups
color-coded.

\begin{figure}[t]
  \centering
  \anvilfig{fig1-geometry.pdf}{fig1\_geometry.py}
  \caption{Benchmark geometry: the DeviceLayout.jl v1.15.0
  \texttt{SingleTransmon} transmon and readout resonator, with the seven
  named physical groups color-coded. Both solvers consume the identical
  mesh generated from this geometry (22{,}684 nodes; 133{,}314
  tetrahedra).}
  \label{fig:geometry}
\end{figure}

\subsection{The junction model: a lumped reactive shunt}
\label{sec:junction}

The Josephson junction is modeled identically in both solvers as a lumped
reactive shunt --- a parallel $L \parallel C$ element with
$L = 14.860$\,nH and $C = 5.5$\,fF --- attached to the four-triangle
\texttt{lumped\_element} surface group $\Gamma$. Following the uniform
lumped-port construction, with port length $\ell$ and width $w$, the shunt
contributes
\begin{equation}
  K_{\mathrm{port}} = \frac{\ell}{w\,L}\, S_\Gamma
  \qquad\text{and}\qquad
  M_{\mathrm{port}} = \frac{C\,\ell}{w}\, S_\Gamma ,
  \label{eq:port}
\end{equation}
where $S_\Gamma$ is the port-surface matrix assembled on $\Gamma$. The
inductive term is frequency-independent and adds to the stiffness matrix;
the capacitive term adds to the mass matrix. The eigenproblem is then the
real symmetric generalized pencil
\begin{equation}
  \left(K + K_{\mathrm{port}}\right) x
  \;=\; \omega^2 \left(M + M_{\mathrm{port}}\right) x .
  \label{eq:pencil}
\end{equation}
In Palace, the same physics is entered as a reactive \texttt{LumpedPort}
($L \parallel C$, no resistance) on the same surface group; the readout
ports are configured open (no $50\,\Omega$ termination), so Palace's pencil
is likewise real. This one-to-one correspondence of the junction treatment
is what makes the comparison same-physics as well as same-mesh.

\section{Solver Architectures Under Comparison}
\label{sec:solvers}

The two solvers share the mathematical problem
(Eq.~\ref{eq:pencil}) and nothing else: no code, no language, no linear
algebra, no eigensolver strategy. Table~\ref{tab:architectures} summarizes
the contrast.

\begin{table}[t]
  \centering
  \caption{Solver architectures compared. The two implementations share the
  discretized problem definition (identical mesh, first-order
  N\'ed\'elec elements, identical lumped-shunt junction) and are otherwise
  independent.}
  \label{tab:architectures}
  \footnotesize
  \setlength{\tabcolsep}{4pt}
  \begin{tabular}{@{}lll@{}}
    \toprule
     & \textbf{geode-fem} & \textbf{Palace} \\
    \midrule
    Language / substrate & Rust / Burn tensors (f64) & C++ / MFEM \\
    Assembly & sparse full-tensor N\'ed\'elec & MFEM $H(\mathrm{curl})$ \\
    Eigensolver & shift-invert Lanczos & SLEPc-class Krylov \\
    Inner linear solve & faer sparse LU (direct) & AMS-preconditioned iterative \\
    Spurious-mode control & none (post-hoc filter; \S\ref{sec:spurious-mode}) &
      divergence-free projection \\
    Parallelism & single process & MPI-distributed \\
    \bottomrule
  \end{tabular}
\end{table}

geode-fem assembles the full-tensor N\'ed\'elec system sparsely in double
precision, reduces to interior DOFs, and solves the pencil with a real
shift-invert Lanczos iteration whose inner solves are direct sparse-LU
factorizations. It runs in a single process with no intra-solve parallelism
today. Palace assembles via MFEM, applies a divergence-free projection to
confine the Krylov iteration to the physical subspace, preconditions with
auxiliary-space methods (AMS)
\citep{hiptmair-xu,tzaniovkolevpanayotsvassilevski2018parallel}, and
distributes over MPI ranks. These choices have consequences for both
correctness behavior (Section~\ref{sec:spurious-mode}) and performance
(Section~\ref{sec:perf-cpu}); the point of the benchmark is to expose those
consequences honestly rather than to declare a winner.

\section{Validation Methodology}
\label{sec:methodology}

The benchmark protocol is oracle-first: every claim is anchored either to
an exact analytic quantity or to a committed reference artifact, and the
test suite is designed so that silent pipeline failures are caught by
construction.

\paragraph{Same-mesh, same-physics protocol.}
The mesh is generated once, committed, and pinned by SHA-256. Both solvers
consume it at matched (first) order with the identical junction model
(Section~\ref{sec:junction}). This removes meshing and model entry as
confounders: any residual disagreement is attributable to formulation or
solver implementation.

\paragraph{Analytic tripwires.}
Two exact quantities guard against coincidental or self-consistent-but-wrong
agreement. First, the isolated junction LC resonance
$f_{LC} = 1/(2\pi\sqrt{LC}) = 17.60$\,GHz anchors the expected location of
the junction mode. Second, a Josephson inductance-scaling tripwire: doubling
$L$ must move the junction mode by exactly $1/\sqrt{2}$. In our runs the
junction mode moves from $17.49$ to $12.37$\,GHz, a ratio of $0.7071$ ---
$1/\sqrt{2}$ to the reported precision (Figure~\ref{fig:lscaling}).

\paragraph{Inverse tests that must fail.}
Alongside tests asserting agreement, the suite contains perturbed
configurations (such as the doubled-$L$ case) that are \emph{required to
disagree} with the committed oracle. A pipeline that still ``passes'' under
such a perturbation is not measuring anything; the inverse tests make that
failure mode loud.

\paragraph{Participation-based mode identification.}
Modes are identified, not eyeballed. For each eigenvector $x$ the junction
participation
\begin{equation}
  p \;=\; \frac{x^{\mathsf{T}} K_{\mathrm{port}}\, x}
               {x^{\mathsf{T}} \left(K + K_{\mathrm{port}}\right) x}
  \label{eq:participation}
\end{equation}
cleanly separates the junction LC mode ($p = 1.000$) from the resonator and
box modes ($p \leq 0.0005$). Palace's per-port energy-participation output
provides the same discriminant on its side, and the two assignments agree.

\paragraph{Determinism check.}
The Palace oracle is not a one-off: rerunning Palace from the committed
configuration reproduces its committed \texttt{eig.csv} bit-for-bit.

\section{Results: Cross-Solver Agreement}
\label{sec:agreement}

Table~\ref{tab:agreement} presents the headline result: on the identical
133k-DOF mesh at matched order, every physical mode agrees to
$\leq 0.03\%$, and the junction LC mode --- the mode the junction model
exists to produce --- agrees to $0.000\%$.

\begin{table}[t]
  \centering
  \caption{Eigenmode agreement on the identical mesh (22{,}684 nodes /
  133{,}314 tets / 156{,}863 N\'ed\'elec DOFs, 133{,}108 interior after
  PEC), first order in both solvers, junction as a lumped reactive shunt
  ($L = 14.860$\,nH, $C = 5.5$\,fF). geode-fem at commit
  \texttt{3174015}; Palace at commit \texttt{fba6a5b}.}
  \label{tab:agreement}
  \small
  \begin{tabular}{@{}lccc@{}}
    \toprule
    Mode & geode-fem (GHz) & Palace (GHz) & $\Delta$ \\
    \midrule
    readout resonator & 5.1528 & 5.1513 & 0.029\% \\
    box/cavity modes & 15.4650 / 18.6927 / 20.6976 / 26.0809 &
      same to $\leq 0.03\%$ & $\leq 0.03\%$ \\
    junction LC & 17.4901 & 17.4901 & 0.000\% \\
    \bottomrule
  \end{tabular}
\end{table}

\begin{figure}[t]
  \centering
  \anvilfig{fig2-agreement.pdf}{fig2\_agreement.py}
  \caption{Mode-frequency agreement: geode-fem versus Palace on the
  identical mesh, with per-mode relative deviations annotated. Data from
  the committed benchmark artifacts
  (\texttt{benchmarks/transmon\_eigen/results.toml} and Palace's committed
  \texttt{eig.csv}).}
  \label{fig:agreement}
\end{figure}

\begin{figure}[t]
  \centering
  \anvilfig{fig3-lscaling.pdf}{fig3\_l\_scaling.py}
  \caption{Josephson inductance-scaling tripwire: junction-mode frequency
  versus junction inductance $L$ on log-log axes. Doubling $L$ moves the
  junction mode from $17.49$ to $12.37$\,GHz --- ratio $0.7071 =
  1/\sqrt{2}$ --- exactly as the $f \propto 1/\sqrt{L}$ scaling requires.}
  \label{fig:lscaling}
\end{figure}

Three features of Table~\ref{tab:agreement} deserve comment. First, the
junction mode sits at $17.4901$\,GHz in both solvers, $0.6\%$ below the
isolated-circuit anchor $f_{LC} = 17.60$\,GHz --- the small shift reflects
the junction's embedding in the field problem, and its near-coincidence
with the analytic value in \emph{both} solvers is the tripwire working as
intended. Second, the mode identification is unambiguous: junction
participation is $p = 1.000$ for the $17.49$\,GHz mode and $p \leq 0.0005$
for every other mode, in both solvers' participation outputs. Third, the
agreement is between genuinely independent code paths --- a direct-factorization
Lanczos in Rust and a projected, preconditioned Krylov method in C++ ---
so a common implementation error is an unlikely explanation for the
residual $\leq 0.03\%$, and the analytic tripwires exclude
self-consistent-but-wrong agreement about the junction physics.

\section{The Absent Qubit Mode: An Expected Negative}
\label{sec:absent-mode}

A reader familiar with the DeviceLayout.jl blog workflow
\citep{devicelayout-blog} will notice that its reported spectrum spans
$[4.14,\ 5.591]$\,GHz --- including a ${\sim}4$\,GHz qubit mode --- while
Table~\ref{tab:agreement} contains no mode below $5.15$\,GHz. This is not a
discrepancy to be minimized; it is a property of the model, and we state it
plainly: \textbf{no ${\sim}4$\,GHz qubit mode exists in this junction model,
by construction.}

The physical transmon frequency arises from the Josephson inductance
resonating against the \emph{pad (shunt) capacitance} of roughly
$80$--$100$\,fF \citep{koch2007}. The v1 lumped-shunt model of
Section~\ref{sec:junction} gives the junction only its own $5.5$\,fF, so
the junction's resonance lands at $17.5$\,GHz, not near $4$\,GHz; the pad
capacitance exists in the field solution but is not wired across the
lumped element. Both solvers agree on the absence: Palace's projected
eigensolve, instructed to hunt at a shift of $\sigma = 4.5$\,GHz, finds
nothing below the $5.15$\,GHz resonator mode. The blog's
$[4.14,\ 5.591]$\,GHz spectrum is therefore \emph{not reproduced by this
model}, and no agreement with it is claimed.

The route to genuine qubit quantities from this benchmark is
energy-participation-ratio (EPR) post-processing over the field modes
\citep{minev2021,nigg2012}, in which the junction's participation in each
electromagnetic mode --- exactly the quantity in
Eq.~\ref{eq:participation} --- feeds the quantization. That analysis is
planned future work (Phase~C of this benchmark program) and is out of scope
for the present cross-validation.

\section{A Disclosed Spurious Mode}
\label{sec:spurious-mode}

Honest cross-validation requires reporting where the solvers
\emph{disagree}, and there is exactly one such place. geode-fem's Lanczos
iteration operates on the raw pencil of Eq.~\ref{eq:pencil} without a
divergence-free projection. On this problem it leaks one spurious mode at
approximately $3.45$\,GHz, below the physical band, localized on the
junction surface with participation $p = 0.994$. Palace, whose eigensolve
projects onto the divergence-free subspace, shows no such mode --- its
spectrum is clean down through the $\sigma = 4.5$\,GHz hunt of
Section~\ref{sec:absent-mode}.

We classify this mode as spurious --- a gradient-nullspace-adjacent
artifact of the un-projected iteration, not a physical resonance --- on
documented criteria: (i) it is absent from Palace's projected spectrum on
the identical mesh; (ii) it sits far below the analytic junction anchor
$f_{LC} = 17.60$\,GHz while carrying near-total junction participation,
a combination inconsistent with the physical junction mode that appears,
in both solvers, at $17.49$\,GHz; and (iii) it does not respond as a
physical mode under the tripwire perturbations. The mode is excluded from
Table~\ref{tab:agreement} by these criteria, which are recorded alongside
the committed benchmark configuration. Figure~\ref{fig:participation}
contrasts the two solvers' participation spectra, spurious mode included.

The fix is known and flagged as future work: a divergence-free or
tree--cotree gauge projection on geode-fem's eigenpath, mirroring the
projection Palace applies. We prefer disclosing the artifact and its
filtering criteria over silently post-selecting the spectrum: the
disagreement is itself a finding about what the divergence-free projection
buys in a production eigensolver.

\begin{figure}[t]
  \centering
  \anvilfig{fig6-participation.pdf}{fig6\_participation.py}
  \caption{Junction-participation spectra, geode-fem versus Palace, on the
  identical mesh. The physical junction LC mode ($p = 1.000$, $17.49$\,GHz)
  appears in both solvers; geode-fem's un-projected Lanczos additionally
  leaks one spurious junction-localized mode near $3.45$\,GHz
  ($p = 0.994$) that Palace's divergence-free projection suppresses.}
  \label{fig:participation}
\end{figure}

\section{Performance: CPU Cell}
\label{sec:perf-cpu}

\subsection{Setup}

All CPU measurements ran on a single Amazon EC2 \texttt{m6i.4xlarge}
instance (8 physical cores / 16 vCPUs, 64\,GB RAM) in
\texttt{us-west-2}.\footnote{On-demand cost approximately \$0.77/hour at
the time of the runs; the entire CPU cell is reproducible for a few
dollars.} Each measurement covers the \emph{full pipeline} --- mesh load,
assembly, solve, and output --- timed with \texttt{/usr/bin/time};
entries with uncertainties are the mean over $n = 3$ runs. geode-fem ran
at commit \texttt{3174015}, Palace at commit \texttt{fba6a5b}, both
release builds.

\subsection{Results}

\begin{table}[t]
  \centering
  \caption{CPU-cell wall-clock and memory, full pipeline (mesh load +
  assembly + solve + output), \texttt{m6i.4xlarge}. Uncertainties are over
  $n = 3$ runs where shown. Palace at 16 ranks (hyperthreads) was refused
  by MPI process binding and is excluded --- it is not a data point.}
  \label{tab:cpu}
  \small
  \begin{tabular}{@{}llcc@{}}
    \toprule
    Solver & Parallelism & Wall (s) & Peak RSS \\
    \midrule
    geode-fem \texttt{@3174015} & 1 process (serial faer sparse-LU
      shift-invert Lanczos) & $51.2 \pm 0.4$ & 3.1\,GB \\
    Palace \texttt{@fba6a5b} & 4 MPI ranks & 50.8 & ${\sim}0.7$\,GB/rank \\
    Palace \texttt{@fba6a5b} & 8 MPI ranks & $30.6 \pm 0.1$ &
      ${\sim}0.5$\,GB/rank \\
    \bottomrule
  \end{tabular}
\end{table}

\begin{figure}[t]
  \centering
  \anvilfig{fig4-cpu-wallclock.pdf}{fig4\_cpu\_wallclock.py}
  \caption{CPU-cell wall-clock comparison (geode-fem single process versus
  Palace at 4 and 8 MPI ranks), with a per-core-efficiency inset. Full
  pipeline timings from Table~\ref{tab:cpu}.}
  \label{fig:cpu}
\end{figure}

Table~\ref{tab:cpu} supports two statements that must be read together.
Per core, geode-fem's serial direct-factorization eigensolve is roughly
$4\times$ more efficient than Palace's distributed Krylov--Schur/AMS solve
on this 133k-DOF problem: one geode-fem core matches four Palace ranks
almost exactly ($51.2$ versus $50.8$\,s). At the whole-box level, Palace's
MPI parallelism wins $1.7\times$ ($30.6$ versus $51.2$\,s), because
geode-fem has no intra-solve parallelism today and leaves seven cores
idle.

This is an architecture trade-off, not a ranking. At 133k DOFs the sparse
LU factorization fits comfortably in memory (3.1\,GB peak) and its flop
efficiency dominates; a distributed iterative method carries communication
and preconditioner-setup overheads that cannot amortize at this size, but
distributes memory (${\sim}0.5$\,GB/rank at 8 ranks) and scales to problem
sizes and rank counts where a serial direct factorization is not viable at
all. The honest summary is that each architecture wins the regime it was
designed for, and this benchmark documents where the crossover sits for a
production transmon geometry at first order. Palace at 16 ranks on the
hyperthreaded vCPUs was refused by MPI process binding; we exclude it
rather than report a degraded configuration as if it were a measurement.

\section{Performance: GPU Cell (Pending)}
\label{sec:perf-gpu}

\TBDGPU{This entire section awaits measured numbers. No GPU timing or
correctness figure below is a result; every value is a placeholder.}

\subsection{Scope and honest framing}

The GPU cell is scoped but not yet run; we state its design and its
limitations now so the eventual numbers cannot be over-read. geode-fem's
GPU path today accelerates the \emph{driven} (deterministic
frequency-sweep) solve --- a matrix-free matvec with an on-device COCG
iteration --- and \emph{not} the eigensolve that produced the headline
comparison of Section~\ref{sec:agreement}. A GPU column for the eigenmode
benchmark would therefore be misleading; the GPU cell instead measures the
driven-solve pipeline on the same fixture. Additionally, the CUDA backend
(burn-cuda 0.21) is f32-only --- the underlying cubecl runtime currently
disables f64\footnote{Upstream constraint of the burn/cubecl stack at
version 0.21 \citep{burn}; tracking-issue URL to be added at revision.
TODO(operator).} --- so all GPU results are mixed-precision-qualified and
will be reported with explicit correctness deltas against the f64 CPU
path, not as drop-in equivalents.

\subsection{Planned measurement}

Hardware: a single Amazon EC2 \texttt{g6e.xlarge} instance (1$\times$
NVIDIA L40S, 46\,GB; CUDA 13.2). Pending at the time of drafting:
\begin{itemize}
  \item geode-fem CUDA-f32 correctness smokes (matrix-free matvec,
    on-device COCG, driven-solve wiring): \TBDGPU{pass/fail + tolerance
    table}.
  \item Driven-solve wall-clock, CPU versus GPU, same fixture:
    \TBDGPU{timing table}.
  \item Palace-GPU (libCEED/CUDA build): \TBDGPU{explicitly future work if
    not run for the revision --- will be stated as such, not silently
    dropped}.
\end{itemize}

% Figure plan slot 5 (GPU cell results) is intentionally not a figure
% environment yet: no data exists. pub-figures must not render a fig5
% until the GPU cell lands. \TBDGPU marker stands in its place:
\TBDGPU{Figure 5 (GPU-cell results) is deferred until measured data
exists.}

\section{Reproducibility}
\label{sec:repro}

Every artifact behind every number in this paper is committed and
scripted.\footnote{Repository URL and archival DOI for the geode-fem
source at commit \texttt{3174015} to be added at submission.
TODO(operator).}

\begin{itemize}
  \item \textbf{Geometry and mesh.} Generated by DeviceLayout.jl v1.15.0's
    \texttt{SingleTransmon} example \citep{devicelayout}; the fixture
    generator is committed, and the MSH~4.1 fixture is pinned by SHA-256
    (\texttt{5b3ff4c3\ldots{}b33dd}; full hash and provenance in the
    committed \texttt{transmon\_smoke.provenance.txt}).
  \item \textbf{The MSH conversion gotcha.} MFEM's MSH~4.1 reader rejects
    the DeviceLayout multi-entity-block file (``vertices indices are not
    unique''), so the Palace runs consume an MSH~2.2 conversion
    (\texttt{gmsh -save -format msh2}). The conversion is node-preserving
    (all 22{,}684 nodes unchanged, physical groups preserved) and
    physics-neutral: eigenvalues reproduce bit-for-bit across the
    conversion. We document this because it is exactly the kind of silent
    pipeline step a same-mesh claim must disclose.
  \item \textbf{Solver configurations.} The Palace configuration and its
    provenance file are committed under
    \path{reference/fixtures/transmon_palace/}; geode-fem's benchmark
    definitions live in \path{benchmarks/transmon_eigen/results.toml},
    alongside the exact benchmark commands.
  \item \textbf{Oracles.} Palace's raw outputs (\texttt{eig.csv},
    per-port participation, run log) are committed; a rerun from the
    committed configuration reproduces \texttt{eig.csv} bit-for-bit.
  \item \textbf{Hardware disclosure.} CPU cell: EC2 \texttt{m6i.4xlarge}
    (8 physical cores / 16 vCPU, 64\,GB), \texttt{us-west-2}. GPU cell
    (pending): EC2 \texttt{g6e.xlarge} (1$\times$ NVIDIA L40S 46\,GB,
    CUDA 13.2).
  \item \textbf{Versions.} geode-fem commit \texttt{3174015}; Palace
    commit \texttt{fba6a5b}; DeviceLayout.jl v1.15.0; Burn/burn-cuda 0.21.
\end{itemize}

\section{Discussion}
\label{sec:discussion}

\paragraph{What the agreement does and does not establish.}
In V\&V terms \citep{roache1997quantification,oberkampf2010verification},
this study is code verification and code-to-code comparison, not
validation against nature: two independent implementations agreeing to
$\leq 0.03\%$ on the identical discrete problem, with analytic tripwires
pinning the junction physics. It does not certify that either solver's
answer matches a measured device --- that is what
\citet{ye2025electromagnetic}-style experimental benchmarks address, and
the two evidence types are complementary. What same-mesh agreement
uniquely rules out is the large class of errors that live in formulation,
assembly, boundary-condition handling, and eigensolver machinery.

\paragraph{Threats to validity.}
A common-mode error --- both solvers wrong in the same way --- cannot be
excluded by cross-comparison alone; the analytic anchors ($f_{LC}$, the
$1/\sqrt{2}$ scaling) and the genuinely disjoint implementations make it
unlikely for the junction physics specifically. Configuration-entry
mismatch (comparing subtly different problems) is mitigated by committed
configurations, the shared mesh file, and participation-based mode
identification on both sides. Performance numbers are $n = 3$ full-pipeline
measurements on one instance type; they support the architecture-level
reading in Section~\ref{sec:perf-cpu}, not microsecond-level claims.

\paragraph{Limitations.}
(i) One geometry, one mesh, first order only --- higher-order and
mesh-refinement convergence is future work (a $p = 2$ Palace
configuration is already committed for that purpose). (ii) The lossless v1
model leaves readout ports open; resistive terminations (and hence
complex-symmetric pencils and mode $Q$s) are out of scope. (iii) The
spurious-mode filter of Section~\ref{sec:spurious-mode} currently leans on
the committed oracle and documented criteria; the principled fix
(divergence-free/tree--cotree projection) is identified but not yet
implemented. (iv) The per-core efficiency result is specific to the
133k-DOF regime where a serial direct factorization is memory-feasible;
it will not extrapolate to problem sizes that demand distribution.
(v) The GPU cell is pending; no GPU claims are made in this version.

\section{Conclusion}
\label{sec:conclusion}

We presented the first independent, quantitative validation of the
Palace and DeviceLayout.jl superconducting-transmon simulation stack: an
independently implemented Rust FEM solver reproduces Palace's eigenmodes on
the identical mesh and identical junction model to $\leq 0.03\%$ across
all physical modes ($0.000\%$ on the junction LC mode), with exact analytic
tripwires guarding the junction physics. The benchmark's methodology ---
same-mesh protocol, oracle-first tripwires, inverse tests that must fail,
participation-based mode identification, and honest-negative reporting of
both the by-construction absent qubit mode and a disclosed spurious mode
--- is itself a reusable template for cross-solver validation in this
domain. The CPU performance cell documents a clean architecture trade-off:
${\sim}4\times$ per-core efficiency for the serial direct-factorization
eigensolve against a $1.7\times$ whole-box win for MPI distribution at
133k DOFs. A mixed-precision GPU cell for the driven solve is scoped and
pending. Both solvers emerge with their correctness confirmed on this
problem class --- which is precisely the outcome an ecosystem that
increasingly trusts open solvers for quantum-hardware design should demand,
and until now could not cite.

\paragraph{Acknowledgments.}
geode-fem's development was AI-orchestrated (agent-driven implementation
with human direction and review); this benchmark program and paper follow
the same discipline, with every reported number traced to a committed,
human-auditable artifact.
% TODO(operator): finalize acknowledgment wording and author list before
% submission.

\bibliographystyle{plainnat}
\bibliography{refs}

\end{document}
